\documentclass[12pt,a4paper]{article}
% kompilirarai s - xelatex -interaction=nonstopmode modelling.tex

\usepackage{fontspec}
\setmainfont{Times New Roman}
\setsansfont{Arial}
\setmonofont{Consolas}

\usepackage{amsmath,amssymb,bm}

\usepackage[margin=2.5cm]{geometry}
\setlength{\parskip}{6pt}
\setlength{\parindent}{0pt}
\usepackage{booktabs}

\usepackage[colorlinks=true,linkcolor=blue,citecolor=blue]{hyperref}

\title{\textbf{SALUS -- Математически модел (само deep learning детектори)}}
\author{FMI\{CODES\} Demo}
\date{\today}

\begin{document}
\maketitle
\tableofcontents
\newpage

\section{Въведение}

Системата за оценяване на изображения използва 3 модела:
\begin{itemize}
  \item Deepfake детектор (ResNet-18, бинарна класификация);
  \item NSFW детектор (Hugging Face image-classification pipeline);
  \item Flux детектор (SigLIP image classifier).
\end{itemize}

\section{Нотация}

Нека входното изображение е $I \in \mathbb{R}^{H \times W \times 3}$,
а множеството достъпни разпознаватели е
\begin{equation}
\mathcal{R} = \{\text{deepfake},\; \text{nsfw},\; \text{flux}\}.
\end{equation}

За всеки разпознавател $r \in \mathcal{R}$ дефинираме:
\begin{equation}
\text{result}_r = (y_r,\; \ell_r,\; s_r,\; \mathbf{p}_r),
\end{equation}
където:
\begin{itemize}
  \item $y_r \in \{0,1\}$ е бинарното решение \emph{is\_target};
  \item $\ell_r$ е текстов етикет (label);
  \item $s_r \in [0,1]$ е основният score;
  \item $\mathbf{p}_r$ е списък от всички предсказания (label, score).
\end{itemize}

\section{Deepfake модел (ResNet-18)}

\subsection{Лицево изрязване (face crop)}

Нека детектираният bbox на най-голямото лице е $(x,y,w,h)$ и padding $\alpha=0.15$.
Крайният crop е:
\begin{align}
x_1 &= \max(0,\; x-\alpha w), & x_2 &= \min(W,\; x+w+\alpha w), \\
y_1 &= \max(0,\; y-\alpha h), & y_2 &= \min(H,\; y+h+\alpha h).
\end{align}

\subsection{Мрежа и решение}

Остатъчен блок в ResNet има вида:
\begin{equation}
\mathbf{z}_{k+1} = \mathcal{F}(\mathbf{z}_k;\theta_k) + \mathbf{z}_k.
\end{equation}

След resize и normalize към $224\times224$, мрежата връща логит $u_{df}$,
а вероятността за deepfake е:
\begin{equation}
p_{df} = \sigma(u_{df}) = \frac{1}{1+e^{-u_{df}}}.
\end{equation}

С праг $\tau_{df}$:
\begin{equation}
y_{df} = \mathbf{1}[p_{df} \ge \tau_{df}], \qquad
\ell_{df} =
\begin{cases}
\text{Deepfake}, & y_{df}=1,\\
\text{Real}, & y_{df}=0,
\end{cases}
\qquad s_{df}=p_{df}.
\end{equation}

\section{Как обучихме нашия Deepfake модел}

\subsection{Данни и етикети}

Моделът е трениран върху датасета
\texttt{muhammadbilal6305/200k-real-vs-ai-visuals-by-mbilal},
като задачата е бинарна класификация:
\begin{equation}
\mathcal{D} = \{(x_i, y_i)\}_{i=1}^{N}, \qquad y_i \in \{0,1\},
\end{equation}
където $y_i=0$ означава \textit{real}, а $y_i=1$ означава \textit{fake}.

\subsection{Разделяне train/validation}

Използваме стратифицирано разделяне (20\% validation):
\begin{equation}
|\mathcal{D}_{train}| \approx 0.8N, \qquad |\mathcal{D}_{val}| \approx 0.2N,
\end{equation}
като относителният дял на класовете се запазва във всеки split.

\subsection{Предобработка и аугментации}

За train използваме resize, random resized crop, horizontal flip,
color jitter и нормализация по канали; за validation -- детерминистичен
resize и същата нормализация.

Нормализацията на пикселите е:
\begin{equation}
\tilde{x}_{c,h,w} = \frac{x_{c,h,w} - \mu_c}{\sigma_c},
\end{equation}
с
\begin{equation}
\bm{\mu}=(0.485,\,0.456,\,0.406),
\qquad
\bm{\sigma}=(0.229,\,0.224,\,0.225).
\end{equation}

\subsection{Балансиране на класовете}

Използваме два механизма за баланс:
\begin{enumerate}
  \item \textbf{WeightedRandomSampler} с тегло на проба
  \begin{equation}
  w_i = \frac{1}{n_{y_i}},
  \end{equation}
  където $n_{y_i}$ е броят примери от класа на $x_i$;
  \item \textbf{BCEWithLogitsLoss} с
  \begin{equation}
  \text{pos\_weight} = \frac{n_{real}}{n_{fake}}.
  \end{equation}
\end{enumerate}

\subsection{Функция на загуба и оптимизация}

Нека $z_i$ е логитът за проба $i$, $p_i=\sigma(z_i)$.
Тренираме с претеглена бинарна крос-ентропия:
\begin{equation}
\mathcal{L}_{BCE} = -\frac{1}{M}\sum_{i=1}^{M}
\left[\alpha y_i \log p_i + (1-y_i)\log(1-p_i)\right],
\end{equation}
където $\alpha=\text{pos\_weight}$ и $M$ е размерът на batch-а.

Оптимизацията е с AdamW:
\begin{equation}
\theta \leftarrow \theta - \eta\,\nabla_{\theta}\mathcal{L}_{BCE},
\end{equation}
с learning rate $\eta=3\times10^{-4}$ и weight decay $10^{-4}$.

\subsection{Оценка и избор на checkpoint}

На validation сета използваме праг $0.5$:
\begin{equation}
\hat{y}_i = \mathbf{1}[p_i \ge 0.5].
\end{equation}

След всяка епоха смятаме Accuracy, Precision, Recall, $F_1$ и ROC-AUC,
като запазваме \textit{best checkpoint} по максимална стойност на
$F_1$ върху validation.

\section{NSFW модел}

NSFW разпознавателят връща вероятности за класове
$\mathcal{C}_{nsfw}=\{\text{normal},\text{nsfw}\}$.

Нека $p_{nsfw}=P(c=\text{nsfw}\mid I)$. Тогава:
\begin{equation}
y_{nsfw}=\mathbf{1}[p_{nsfw}\ge\tau_{nsfw}], \qquad s_{nsfw}=p_{nsfw}.
\end{equation}

Етикетът е дефиниран така, че при нисък NSFW score да връща normal:
\begin{equation}
\ell_{nsfw}=
\begin{cases}
\text{nsfw}, & y_{nsfw}=1,\\
\arg\max\limits_{c\in\mathcal{C}_{nsfw}} P(c\mid I), & y_{nsfw}=0.
\end{cases}
\end{equation}

\section{Flux модел (SigLIP)}

Моделът връща logits $\mathbf{z}\in\mathbb{R}^{K}$ и Softmax вероятности:
\begin{equation}
\mathbf{p}_{flux} = \text{softmax}(\mathbf{z}),
\qquad
p_{flux} = P(c=\text{Flux.1\_Generated}\mid I).
\end{equation}

С праг $\tau_{flux}$:
\begin{equation}
y_{flux}=\mathbf{1}[p_{flux}\ge\tau_{flux}],\qquad s_{flux}=p_{flux},
\end{equation}
\begin{equation}
\ell_{flux}=
\begin{cases}
\text{Flux.1\_Generated}, & y_{flux}=1,\\
\arg\max\limits_{c} P(c\mid I), & y_{flux}=0.
\end{cases}
\end{equation}

\section{Multi-endpoint формулировка}

За подмножество заявени разпознаватели $\mathcal{R}_{req}\subseteq\mathcal{R}$,
endpoint-ът изпълнява всеки модел паралелно и връща:
\begin{equation}
\text{results} = \{r \mapsto \text{result}_r \mid r\in\mathcal{R}_{req}\}.
\end{equation}

Изображението се идентифицира с SHA-256 хеш:
\begin{equation}
h = \text{SHA256}(\text{bytes}(I)).
\end{equation}

Практическо агрегирано правило за флагване (ако е нужно на продуктово ниво):
\begin{equation}
y_{global} = \max_{r\in\mathcal{R}_{req}} y_r,
\end{equation}
т.е. изображението се счита за рисково, ако поне един разпознавател го маркира.

\section{Метрики за оценка}

За всеки модел поотделно се използват:
\begin{align}
\text{Precision} &= \frac{TP}{TP+FP+\varepsilon}, \\
\text{Recall}    &= \frac{TP}{TP+FN+\varepsilon}, \\
F_1              &= \frac{2\,\text{Precision}\cdot\text{Recall}}
                         {\text{Precision}+\text{Recall}+\varepsilon}, \\
\text{Accuracy}  &= \frac{TP+TN}{TP+TN+FP+FN+\varepsilon}.
\end{align}

\section{Кратко обобщение}

\begin{center}
\begin{tabular}{lll}
\toprule
\textbf{Модел} & \textbf{Основен score} & \textbf{Праг / решение} \\
\midrule
Deepfake (ResNet-18) & $s_{df}=p_{df}$ & $y_{df}=\mathbf{1}[p_{df}\ge\tau_{df}]$ \\
NSFW pipeline & $s_{nsfw}=p_{nsfw}$ & $y_{nsfw}=\mathbf{1}[p_{nsfw}\ge\tau_{nsfw}]$ \\
Flux (SigLIP) & $s_{flux}=p_{flux}$ & $y_{flux}=\mathbf{1}[p_{flux}\ge\tau_{flux}]$ \\
\bottomrule
\end{tabular}
\end{center}

\end{document}
